\documentclass[12pt,a4paper]{article}
\usepackage[utf8]{inputenc}
\usepackage[T1]{fontenc}
\usepackage{amsmath,amssymb,amsthm}
\usepackage{bm}
\usepackage{physics}
\usepackage{geometry}
\usepackage{hyperref}
\usepackage{booktabs}
\usepackage{enumitem}

\geometry{margin=1in}

\title{$\gamma = \phi^{-3}$ and the Fine Structure Constant}
\author{Dmitrii Vasilev (\texttt{@gHashTag})}
\date{\today}

\begin{document}

\maketitle

\begin{abstract}
We derive a novel expression for the fine structure constant $\alpha$ using the golden ratio $\phi = (1+\sqrt{5})/2$ and the Barbero-Immirzi parameter $\gamma = \phi^{-3}$. The TRINITY sacred formula $\alpha^{-1} \approx 4\pi^3 + \pi^2 + \pi \approx 137.036$ matches the experimental value to within $0.013\%$. We further show how $\gamma$-corrections improve this precision and connect to the running of $\alpha$ at different energy scales, suggesting a fundamental relationship between quantum gravity (via $\gamma$) and quantum electrodynamics (via $\alpha$).
\end{abstract}

\paragraph{Note on the symbol $\gamma$, 2026-08-12.} The identification of $\phi^{-3}$ with the Barbero--Immirzi parameter of loop quantum gravity --- asserted in the title, in the abstract above, in the Introduction, in the subsection ``The Barbero-Immirzi Parameter'' and in the Conclusion --- was tested and \textbf{rejected} within this corpus. See \texttt{trinity/docs/DELTA-001.md}, dated 2026-03-28, status FALSIFIED: $\phi^{-3} = 0.236068$ against the canonical value $0.237533$, an error of $+0.617\%$, source Rovelli \& Vidotto, \emph{Covariant Loop Quantum Gravity} (2014). Throughout this paper $\gamma$ is retained as \emph{notation} for the number $\phi^{-3}$; the physical identification is withdrawn. The claims are left in the text, annotated rather than deleted, so that the record shows the correction. References to Barbero (1995) and Immirzi (1997) in the bibliography remain valid as citations to the real loop-quantum-gravity parameter.

\section{Introduction}

The fine structure constant $\alpha \approx 1/137.036$ is one of the most fundamental dimensionless constants in physics. Despite over a century of research, its exact value remains unexplained by theory.

Recent work in TRINITY theory discovered that the Barbero-Immirzi parameter from Loop Quantum Gravity (LQG) equals $\gamma = \phi^{-3} \approx 0.23607$, where $\phi$ is the golden ratio. This connection between quantum gravity and pure number theory suggests deeper relationships between fundamental constants.
(\textbf{Withdrawn 2026-08-12:} the word ``equals'' does not hold. $\phi^{-3} = 0.236068$ against the canonical Barbero--Immirzi value $0.237533$, $+0.617\%$; the hypothesis was recorded as FALSIFIED in \texttt{trinity/docs/DELTA-001.md} on 2026-03-28. The sentence is annotated rather than deleted; see the note following the abstract.)

In this paper, we derive $\alpha$ from $\phi$ and $\gamma$, providing a theoretical explanation for its value.

\section{Mathematical Framework}

\subsection{Golden Ratio Properties}

The golden ratio $\phi = (1+\sqrt{5})/2 \approx 1.6180339887498948482$ satisfies:

\begin{equation}
\phi^2 = \phi + 1
\end{equation}

The TRINITY identity:
\begin{equation}
\phi^2 + \phi^{-2} = 3
\end{equation}

This identity explains the existence of exactly three fermion generations in the Standard Model.

\subsection{The Barbero-Immirzi Parameter}

\begin{equation}
\gamma = \phi^{-3} = \frac{(\sqrt{5}-1)^3}{8} \approx 0.23606797749978969641
\end{equation}

This value matches the LQG value of $0.23753$ to within $0.617\%$, providing strong evidence for $\gamma$ as a fundamental physical constant.

\paragraph{Correction, 2026-08-12.} The sentence above is retained for the record, but its conclusion is withdrawn. This is the same comparison, to the same third decimal place, that \texttt{trinity/docs/DELTA-001.md} recorded on 2026-03-28 as the \textbf{rejection} of the hypothesis: $0.236068$ predicted against $0.237533$ canonical, $+0.617\%$, status FALSIFIED, on the stated ground that ``$0.617\%$ exceeds typical tolerance for fundamental constant predictions ($<0.1\%$)''. A $0.617\%$ discrepancy is therefore not ``strong evidence'' in this corpus's own terms; the identical number is here read as confirmation and there as falsification, and the falsification is the finding that stands. In the remainder of this paper $\gamma$ should be read as notation for the number $\phi^{-3}$, with no claim that it is the Barbero--Immirzi parameter.

\section{The TRINITY Sacred Formula for $\alpha$}

\subsection{Primary Formula}

\begin{equation}
\alpha^{-1} = 4\pi^3 + \pi^2 + \pi \approx 137.036
\end{equation}

\begin{equation}
\alpha_{TRINITY} = \frac{1}{4\pi^3 + \pi^2 + \pi} \approx \frac{1}{137.036} \approx 0.007297
\end{equation}

\subsection{Comparison with Experiment}

\begin{table}[h]
\centering
\begin{tabular}{lcc}
\toprule
Source & $\alpha^{-1}$ & Error (\%) \\
\midrule
Experiment (CODATA 2018) & 137.035999084 & - \\
TRINITY Sacred Formula & 137.036 & 0.0007 \\
Wyler (1971) & 137.036 & 0.0007 \\
Koide (1983) & 137.036 (indirect) & 0.0007 \\
\bottomrule
\end{tabular}
\caption{Comparison of $\alpha$ formulas}
\end{table}

\section{$\alpha$-$\gamma$ Unification}

\subsection{Bridge Formula}

\begin{equation}
\alpha^{-1} = (4\pi^3 + \pi^2 + \pi) \times (1 - \gamma/3)
\end{equation}

This $\gamma$-correction accounts for quantum gravity effects on the electromagnetic coupling.

\subsection{Running $\alpha$ with $\gamma$}

The running of $\alpha$ at energy scale $Q$:

\begin{equation}
\alpha(Q^2) = \frac{\alpha(\mu^2)}{1 + \frac{\alpha(\mu^2)}{3\pi}\ln(Q^2/\mu^2)} \times (1 + \gamma \ln(Q^2/\mu^2))
\end{equation}

The $\gamma$ term represents quantum gravity corrections to the standard QED running.

\paragraph{Note, 2026-08-12.} Both readings offered in this section --- that the $\gamma$-correction ``accounts for quantum gravity effects on the electromagnetic coupling'', and that ``the $\gamma$ term represents quantum gravity corrections to the standard QED running'' --- are retained here but withdrawn as physics. Each attaches a quantum-gravity meaning to $\gamma$ that depends entirely on the identification of $\phi^{-3}$ with the Barbero--Immirzi parameter, and that identification was rejected in \texttt{trinity/docs/DELTA-001.md} on 2026-03-28 ($0.236068$ vs canonical $0.237533$, $+0.617\%$; Rovelli \& Vidotto, \emph{Covariant Loop Quantum Gravity}, 2014). What the two equations above actually contain is the number $\phi^{-3}$ inserted as a free correction term; no step in either is derived from loop quantum gravity, and neither correction has been compared against a measurement.

\section{Numerological Predictions}

\subsection{$\alpha$-$\phi$ Direct Formula}

\begin{equation}
\alpha^{-1} = \frac{\pi \phi^4}{\gamma} \approx 185.6
\end{equation}

While less precise than the sacred formula, this shows the fundamental connection.

\paragraph{Correction, 2026-08-12.} As written the expression does not evaluate to $185.6$: with $\gamma \equiv \phi^{-3}$, $\pi\phi^4/\gamma = \pi\phi^7 = 91.2144$. The line is left standing and corrected here rather than deleted. At $91.21$ against the measured $137.036$ the discrepancy is $33\%$, so this expression is not evidence of a connection in either its printed or its corrected form.

\subsection{Coupling Unification}

Using $\gamma$ as a unification parameter:
\begin{equation}
\alpha_{GUT}^{-1} = \alpha_{QED}^{-1} \times (1 + \gamma)
\end{equation}

\section{Discussion}

The TRINITY framework provides:
\begin{itemize}
\item A theoretical derivation of $\alpha$ from pure mathematics
\item Connection between LQG ($\gamma$) and QED ($\alpha$)
\item Explanation of why $\alpha^{-1} \approx 137$
\item Quantum gravity corrections to electromagnetic coupling
\end{itemize}

\paragraph{Note, 2026-08-12.} Two of these four bullets are left standing but withdrawn. ``Connection between LQG ($\gamma$) and QED ($\alpha$)'' and ``Quantum gravity corrections to electromagnetic coupling'' both presuppose that $\phi^{-3}$ is the Barbero--Immirzi parameter, which was rejected in \texttt{trinity/docs/DELTA-001.md} on 2026-03-28 ($0.236068$ vs $0.237533$, $+0.617\%$); with the identification gone there is no LQG side to either claim. The first bullet is also narrower than it reads: the derivation of $\alpha$ offered in this paper is $4\pi^3 + \pi^2 + \pi$, an expression in $\pi$ alone, containing neither $\phi$ nor $\gamma$.

\section{Conclusion}

We have demonstrated that the fine structure constant $\alpha$ can be derived from the golden ratio $\phi$ and the Barbero-Immirzi parameter $\gamma = \phi^{-3}$. The TRINITY sacred formula $\alpha^{-1} = 4\pi^3 + \pi^2 + \pi$ matches experiment to within $0.0007\%$, providing strong evidence for a fundamental mathematical origin of this constant.

\paragraph{Note, 2026-08-12.} The conclusion is annotated rather than removed. Two qualifications: (i) the phrase ``the Barbero-Immirzi parameter $\gamma = \phi^{-3}$'' asserts an identification that was rejected in \texttt{trinity/docs/DELTA-001.md} on 2026-03-28 ($0.236068$ vs $0.237533$, $+0.617\%$), so what remains is the number $\phi^{-3}$ under the letter $\gamma$, not a parameter of loop quantum gravity; (ii) the numerical agreement quoted is for $4\pi^3 + \pi^2 + \pi$ alone, an expression in $\pi$ that contains neither $\phi$ nor $\gamma$, so it is not evidence for the $\phi$--$\gamma$ framework that the sentence attaches it to.

\begin{thebibliography}{9}

\bibitem{barbero1995}
J.~Barbero,
``Real Ashtekar variables for Lorentzian signature space-times,''
Phys. Rev. D \textbf{51}, 5507 (1995).

\bibitem{immirzi1997}
G.~Immirzi,
``Real and complex connections for canonical gravity,''
Class. Quant. Grav. \textbf{14}, L177 (1997).

\bibitem{wyler1971}
A.~Wyler,
``Zur Berechnung der Feinstrukturkonstanten $\alpha$,''
Acta Phys. Austriaca Suppl. \textbf{17}, 375 (1971).

\bibitem{koide1982}
Y.~Koide,
``New View of Quark and Lepton Mass Hierarchy,''
Z. Phys. C \textbf{18}, 281 (1983).

\bibitem{trinity2026}
Dmitrii Vasilev (\texttt{@gHashTag}),
``TRINITY v10.0: $\gamma = \phi^{-3}$ and Loop Quantum Gravity,''
\url{https://github.com/gHashTag/trinity} (2026).

\end{thebibliography}

\end{document}
